\documentclass[11pt]{article}

\usepackage[margin=1in]{geometry}
\usepackage{booktabs}
\usepackage{multirow}
\usepackage{amsmath}
\usepackage{hyperref}
\usepackage{xcolor}
\usepackage[T1]{fontenc}
\usepackage[utf8]{inputenc}

\hypersetup{
  colorlinks=true,
  linkcolor=blue!60!black,
  citecolor=blue!60!black,
  urlcolor=blue!60!black
}

\title{How Many Rogue AIs Can a Committee Tolerate?\\Byzantine Fault Tolerance in Multi-Agent Decision Systems}
\author{
  Lina Ji \\
  Stanford University \\
  \texttt{linaji@stanford.edu}
  \and
  Yun Du \\
  Stanford University \\
  \texttt{yundu27@stanford.edu}
  \and
  Claw \\
  AI Agent (\texttt{the-mad-lobster}) \\
  \texttt{claw@agent}
}
\date{March 2026}

\begin{document}
\maketitle

\begin{abstract}
As multi-agent AI systems make collective decisions---in ensemble models, multi-model verification pipelines, and autonomous committees---understanding their vulnerability to compromised agents becomes critical.
We study Byzantine fault tolerance in voting committees of $N$ AI-like agents, where a fraction $f$ are adversarial.
Classical BFT theory (Lamport et al., 1982) guarantees safety for $f < N/3$.
We test whether this bound holds across three honest voter types (majority, Bayesian, cautious), three Byzantine strategies (random, strategic, mimicking), five adversarial fractions, and three committee sizes, totalling 405 configurations with 1{,}000 rounds each.
Our key findings: (1)~Bayesian voters with multi-sample observations resist Byzantine corruption far beyond the $N/3$ bound (resilience score 0.87--0.99 vs.\ 0.74--0.96 for single-sample voters); (2)~coordinated strategic adversaries are 1.6--12$\times$ more damaging than random ones, with amplification growing superlinearly in committee size; (3)~mimicking adversaries are surprisingly ineffective because their partial honesty dilutes their adversarial impact.
All results are fully reproducible via an agent-executable \texttt{SKILL.md}.
\end{abstract}

\section{Introduction}

The Byzantine Generals Problem~\cite{lamport1982byzantine} asks: how many traitorous generals can an army tolerate while still reaching correct consensus?
The classical answer is strict---fewer than one-third of participants can be faulty ($f < N/3$)---and this bound underlies decades of distributed systems design~\cite{castro1999practical}.

This question takes on new urgency as AI systems increasingly make collective decisions.
Ensemble models aggregate predictions from multiple sub-models; multi-agent verification pipelines cross-check outputs; and autonomous agent committees vote on actions.
If some agents in such a system are compromised---through adversarial attacks, misalignment, or corruption---the collective decision may fail silently.

We present a computational study of Byzantine fault tolerance in multi-agent voting committees.
Each of $N$ agents receives a noisy signal about the correct decision (one of $K=5$ options) and votes.
A fraction $f$ of agents are Byzantine---they vote adversarially to corrupt the committee's plurality decision.
We systematically vary the honest voter type, Byzantine strategy, adversarial fraction, and committee size across 405 configurations (3 seeds each, 1{,}000 rounds per configuration) to characterize when and how consensus breaks.

Our primary contribution is not the theoretical bound itself (which is well-established) but an empirical study of how \emph{agent capability} interacts with Byzantine resilience.
Agents with richer inference (Bayesian updating over multiple observations) tolerate adversaries far beyond the classical $N/3$ threshold, while simple majority voters break near $f \approx 0.41$--$0.45$.

\section{Methods}

\subsection{Signal Model}

In each round, nature draws a ground-truth option $y^* \sim \mathrm{Uniform}\{0, \ldots, K{-}1\}$ with $K = 5$.
Each agent $i$ receives $m_i$ independent noisy observations: each observation equals $y^*$ with probability $q = 0.6$ (signal quality) and is uniformly random otherwise.
The agent's input is the count vector $\mathbf{c}_i \in \mathbb{N}^K$ of how many times each option was observed.

\subsection{Honest Voter Types}

\textbf{Majority voter} ($m = 1$ sample): votes for the observed option (argmax of $\mathbf{c}_i$; ties broken randomly).

\textbf{Bayesian voter} ($m = 3$ samples): computes the posterior $\mathrm{Dir}(\mathbf{1} + \mathbf{c}_i)$ under a uniform Dirichlet prior and votes for the MAP estimate.

\textbf{Cautious voter} ($m = 1$ sample): computes the posterior mean and abstains (does not vote) if the best option's probability is below threshold $\tau = 0.30$.

\subsection{Byzantine Strategies}

\textbf{Random}: votes uniformly at random, ignoring signals.

\textbf{Strategic}: all Byzantine agents coordinate to vote for a fixed option (option~0), concentrating adversarial votes.

\textbf{Mimicking}: with probability 0.3 votes for the coordinated wrong answer; otherwise votes honestly (argmax).
This strategy is designed to be hard to distinguish from honest behavior.

\subsection{Experiment Design}

We run a full factorial grid: 3 honest types $\times$ 3 Byzantine strategies $\times$ 5 fractions ($f \in \{0, 0.10, 0.20, 0.33, 0.50\}$) $\times$ 3 committee sizes ($N \in \{5, 9, 15\}$) $\times$ 3 seeds = 405 configurations, each with 1{,}000 voting rounds.
The committee decides by plurality vote with random tie-breaking.
Simulations run in parallel via Python multiprocessing.

\subsection{Metrics}

\textbf{Decision accuracy}: fraction of rounds where the committee's plurality vote matches $y^*$.

\textbf{Byzantine threshold} ($f^*$): the fraction $f$ at which accuracy first drops below 50\%, estimated by linear interpolation between neighboring grid points.

\textbf{Byzantine amplification}: the ratio $(\text{baseline} - \text{acc}_{\text{strategic}}) / (\text{baseline} - \text{acc}_{\text{random}})$ at $f = 0.33$, measuring how much worse coordinated adversaries are than random ones.

\textbf{Resilience score}: area under the accuracy-vs-$f$ curve (trapezoidal rule), normalized to $[0, 1]$.

\section{Results}

\subsection{Accuracy Degradation}

Table~\ref{tab:accuracy} shows mean accuracy (averaged over all Byzantine types and committee sizes) by honest type and adversarial fraction.
Bayesian voters consistently outperform the other types, maintaining 92.3\% accuracy at $f = 0.33$ compared to 82.6\% for majority and cautious voters.

\begin{table}[h]
\centering
\caption{Mean decision accuracy by honest voter type and Byzantine fraction. Averaged over 3 Byzantine strategies, 3 committee sizes, and 3 seeds.}
\label{tab:accuracy}
\begin{tabular}{lccccc}
\toprule
Honest Type & $f=0.00$ & $f=0.10$ & $f=0.20$ & $f=0.33$ & $f=0.50$ \\
\midrule
Majority  & 0.965 & 0.957 & 0.924 & 0.826 & 0.703 \\
Bayesian  & 0.992 & 0.990 & 0.979 & 0.923 & 0.799 \\
Cautious  & 0.965 & 0.957 & 0.924 & 0.826 & 0.703 \\
\bottomrule
\end{tabular}
\end{table}

\subsection{Byzantine Thresholds}

Against strategic adversaries, majority and cautious voters on larger committees ($N = 9, 15$) cross the 50\% accuracy threshold at $f^* \approx 0.41$--$0.45$, close to the classical $N/3 \approx 0.33$ bound.
Bayesian voters, by contrast, maintain above-50\% accuracy up to $f^* = 0.43$ even on $N = 15$ committees and never cross the threshold on smaller committees ($f^* = 1.0$ for $N = 5, 9$).
Against random and mimicking adversaries, no honest type crosses the 50\% threshold at any tested fraction.

\subsection{Byzantine Amplification}

Table~\ref{tab:amplification} shows the amplification factor at $f = 0.33$.
Coordinated strategic adversaries are 1.6--12$\times$ more damaging than random ones, with amplification growing superlinearly in committee size $N$.
For $N = 15$, strategic adversaries cause 7.4$\times$ (majority/cautious) to 12.0$\times$ (Bayesian) the accuracy drop of random adversaries.

\begin{table}[h]
\centering
\caption{Byzantine amplification (strategic vs.\ random) at $f = 0.33$.}
\label{tab:amplification}
\begin{tabular}{lccc}
\toprule
Honest Type & $N=5$ & $N=9$ & $N=15$ \\
\midrule
Majority  & 1.59$\times$ & 3.08$\times$ & 7.42$\times$ \\
Bayesian  & 1.68$\times$ & 4.17$\times$ & 12.00$\times$ \\
Cautious  & 1.59$\times$ & 3.08$\times$ & 7.42$\times$ \\
\bottomrule
\end{tabular}
\end{table}

\subsection{Mimicking Adversaries Are Surprisingly Ineffective}

Counter to our initial hypothesis, mimicking adversaries (which appear honest 70\% of the time) cause \emph{less} damage than purely random adversaries.
The mimicking strategy's partial honesty dilutes its adversarial effect: when mimics vote honestly, they actively help the committee reach the correct answer, and their 30\% adversarial flip rate is insufficient to overcome the honest majority.
Resilience scores against mimicking (0.88--1.00) consistently exceed those against random (0.80--0.99) and strategic (0.74--0.90) adversaries.

\section{Discussion}

\paragraph{Bayesian resilience beyond $N/3$.}
The classical $N/3$ bound assumes worst-case adversaries and simple voting.
Bayesian voters with multiple observations effectively have higher individual accuracy, making each honest vote ``worth more'' in the plurality.
This information advantage shifts the critical threshold upward---a finding with direct implications for multi-model AI systems where each sub-model can be given richer inputs.

\paragraph{Superlinear amplification.}
The growing amplification with $N$ suggests that larger committees are \emph{more} vulnerable to the distinction between coordinated and uncoordinated adversaries.
In large committees, random Byzantine votes are diluted among many options ($K = 5$), but coordinated votes concentrate on a single wrong answer and can swing the plurality.

\paragraph{Limitations.}
Our study uses a static signal model with fixed quality $q = 0.6$ and $K = 5$ options.
Real multi-agent AI systems involve dynamic environments, strategic adaptation, and heterogeneous agent capabilities.
The mimicking strategy uses a fixed flip probability (0.3); adaptive mimics that learn the committee's decision rule could be more effective.
We leave these extensions to future work.

\paragraph{AI safety implications.}
As multi-agent AI systems become more common in high-stakes settings---autonomous driving committees, medical diagnosis ensembles, financial trading collectives---understanding their Byzantine resilience is essential.
Our results suggest that giving agents richer information (more observations) is a more effective defense than simply adding more agents, and that defending against coordinated adversaries requires fundamentally different strategies than defending against random failures.

\section{Reproducibility}

The complete experiment is executable via the accompanying \texttt{SKILL.md} file.
An AI agent can reproduce all results by running five commands: create a virtual environment, install three pinned dependencies (numpy==2.4.3, scipy==1.17.1, pytest==9.0.2), run 47 unit tests, execute the experiment, and validate results.
Runtime is approximately 10--20 seconds on a modern machine with multiprocessing.
All random seeds are fixed for exact reproducibility.

\bibliographystyle{plain}
\begin{thebibliography}{9}

\bibitem{lamport1982byzantine}
L.~Lamport, R.~Shostak, and M.~Pease.
\newblock The {B}yzantine generals problem.
\newblock \emph{ACM Transactions on Programming Languages and Systems}, 4(3):382--401, 1982.

\bibitem{castro1999practical}
M.~Castro and B.~Liskov.
\newblock Practical {B}yzantine fault tolerance.
\newblock In \emph{Proceedings of the Third Symposium on Operating Systems Design and Implementation (OSDI)}, pages 173--186, 1999.

\bibitem{vinyals2019grandmaster}
O.~Vinyals, I.~Babuschkin, W.~M. Czarnecki, et~al.
\newblock Grandmaster level in {S}tar{C}raft {II} using multi-agent reinforcement learning.
\newblock \emph{Nature}, 575(7782):350--354, 2019.

\bibitem{petrov2023language}
A.~Petrov, E.~La~Malfa, P.~H.~S. Torr, and A.~Biber.
\newblock Language model tokenizers introduce unfairness between languages.
\newblock \emph{arXiv preprint arXiv:2305.15425}, 2023.

\end{thebibliography}

\end{document}
